Theorem 2. Let $I$ be the interior of a polygon in $\mathbf{R}^{2}$, and let $P, Q, R$, and $S$ be points of $\operatorname{Fr} I$, appearing in the stated cyclic order on $\operatorname{Fr} I$. Let $A$ be an arc from $P$ to $R$, lying in $\bar{I}$, such that $A \cap \operatorname{Fr} I=\{P, R\}$. Then $I-A$ is the union of two disjoint open sets $U_{Q}, U_{S}$, containing $Q$ and $S$ in their frontiers.\\
(Note that by Theorem 3.5 there is no loss of generality in supposing that $\bar{I}$ is a rectangular region. Similarly, by the result of Problem 3.5, we can put $P, Q, R$, and $S$ into the positions shown in Figure 4.1.)

Proof. Let $Q^{\prime}$ and $S^{\prime}$ be points of $I$, near $Q$ and $S$, as in the figure. If $Q^{\prime}$ and $S^{\prime}$ are in the same component of $I-A$, then there is a broken line

\begin{figure}[h]
\begin{center}
  \includegraphics[alt={},max width=\textwidth]{c570eb93-f29d-4347-b75d-c251ce2386f6-032_398_697_438_715}
\captionsetup{labelformat=empty}
\caption{Figure 4.1}
\end{center}
\end{figure}

from $Q^{\prime}$ to $S^{\prime}$ in $I-A$. Therefore there is a broken line $B$, from $Q$ to $S$, lying in $\bar{I}-A$ and intersecting $\operatorname{Fr} I$ only at $Q$ and $S$.

But $P$ and $R$ lie in the same component of $\bar{I}-B$, because $A \subset \bar{I}-B$ and $A$ is connected. This contradicts Theorem 2.8. \(\square\)
